% ============================================================
%
%  CIIR STRICT FORMALIZATION
%
%  Rigorous First-Principles Reconstruction of
%  Constraint-Induced Interface Realism
%
%  Tasks:
%    1. Axiomatize Φ with no dual-role ambiguity
%    2. Dimensional consistency audit and correction
%    3. CRS dynamical substrate (explicit update rule)
%    4. Novel falsifiable quantitative prediction
%
% ============================================================

\documentclass[11pt,a4paper]{article}

% ------------------------------------------------------------
%  Packages
% ------------------------------------------------------------
\usepackage[utf8]{inputenc}
\usepackage[T1]{fontenc}
\usepackage{amsmath,amssymb,amsthm}
\usepackage{mathtools}
\usepackage{bm}
\usepackage{hyperref}
\usepackage{cleveref}
\usepackage{enumitem}
\usepackage{booktabs}
\usepackage{array}
\usepackage{geometry}
\usepackage{xcolor}
\geometry{margin=1in}

% ------------------------------------------------------------
%  Theorem-like environments
% ------------------------------------------------------------
\newtheorem{definition}{Definition}[section]
\newtheorem{axiom}[definition]{Axiom}
\newtheorem{theorem}[definition]{Theorem}
\newtheorem{lemma}[definition]{Lemma}
\newtheorem{proposition}[definition]{Proposition}
\newtheorem{corollary}[definition]{Corollary}
\theoremstyle{remark}
\newtheorem{remark}[definition]{Remark}

% ------------------------------------------------------------
%  Custom commands
% ------------------------------------------------------------
\newcommand{\CR}{\mathcal{R}}
\newcommand{\CM}{\mathcal{M}}
\newcommand{\HC}{\mathcal{H}_C}
\newcommand{\HQ}{\mathcal{H}}
\newcommand{\CB}{\mathcal{B}}
\newcommand{\CA}{\mathcal{A}}
\newcommand{\CK}{\mathcal{K}}
\newcommand{\CL}{\mathcal{L}}
\newcommand{\hatC}{\hat{C}}
\newcommand{\hatH}{\hat{H}}
\newcommand{\hatP}{\hat{P}}
\newcommand{\id}{\mathrm{id}}
\newcommand{\Tr}{\mathrm{Tr}}
\newcommand{\ad}{\mathrm{ad}}
\DeclareMathOperator{\spec}{spec}
\DeclareMathOperator{\ran}{ran}
\DeclareMathOperator{\dom}{dom}
\DeclareMathOperator{\rank}{rank}
\DeclareMathOperator{\diag}{diag}
\DeclareMathOperator{\vol}{vol}

% ============================================================
\title{%
  \textbf{Strict Formalization of Constraint-Induced Interface Realism}\\[6pt]
  \large Axiomatization of $\Phi$, Dimensional Audit, CRS Substrate,\\
  and Falsifiable Prediction
}
\author{%
  QRATUM Project\\
  \texttt{robertringler/QRATUM}
}
\date{April 2026}

\begin{document}
\maketitle

\begin{abstract}
We present a rigorous first-principles formalization of the
Constraint-Induced Interface Realism (CIIR) framework.  The document
addresses four deliverables: (1)~a precise axiomatization of the
interface map~$\Phi$ with explicit factorization into compositional
components eliminating dual-role ambiguity, (2)~a complete dimensional
consistency audit of all geometric and informational bounds with
corrected dimensionless forms, (3)~an explicit Computational Reality
System (CRS) dynamical substrate specified as a concrete update rule
sufficient to generate evolution, and (4)~a novel quantitative
prediction---a constraint-curvature-dependent anomalous decoherence
signature---that is falsifiable with near-future quantum-optical
experiments.  We maintain strict separation between formal modeling,
system theory, and speculative extrapolation.  All fail conditions are
explicitly checked.
\end{abstract}

\tableofcontents
\newpage

% %%%%%%%%%%%%%%%%%%%%%%%%%%%%%%%%%%%%%%%%%%%%%%%%%%%%%%%%%%%%
% SECTION 1: FOUNDATIONS AND STATE SPACES
% %%%%%%%%%%%%%%%%%%%%%%%%%%%%%%%%%%%%%%%%%%%%%%%%%%%%%%%%%%%%
\section{Foundations: State Spaces and Measure Structures}
\label{sec:foundations}

We begin by establishing the mathematical infrastructure upon which the
entire formalization rests.  All objects are defined with explicit
domains, codomains, and measure-theoretic properties.

\begin{definition}[Reality Space]
\label{def:reality}
The \emph{reality space} is a triple $(\CR, d_{\CR}, \mathcal{B}(\CR))$
where:
\begin{enumerate}[label=(\roman*)]
  \item $\CR$ is a non-empty set of \emph{ontic states},
  \item $d_{\CR}: \CR \times \CR \to \mathbb{R}_{\geq 0}$ is a
    complete separable metric (making $(\CR, d_{\CR})$ a Polish space),
  \item $\mathcal{B}(\CR)$ is the Borel $\sigma$-algebra generated by
    $d_{\CR}$.
\end{enumerate}
The measurable space $(\CR, \mathcal{B}(\CR))$ is a standard Borel
space.
\end{definition}

\begin{definition}[Constraint Manifold]
\label{def:constraint_manifold}
A \emph{constraint manifold} is a tuple $(\CM, g, \iota, \hatC)$ where:
\begin{enumerate}[label=(\roman*)]
  \item $(\CM, g)$ is a smooth, connected, compact Riemannian manifold
    of dimension $n \geq 1$,
  \item $\iota: \CM \hookrightarrow \CR$ is a $C^2$-smooth isometric
    embedding into the reality space,
  \item $\hatC \in \Gamma(\mathrm{End}(T\CM))$ is a smooth, positive
    semi-definite, bounded endomorphism field (the \emph{constraint
    operator}).
\end{enumerate}
The \emph{constraint curvature} is the scalar field
$\kappa: \CM \to \mathbb{R}_{\geq 0}$ defined by
$\kappa(m) := \|\nabla\hatC(m)\|_g$, where $\nabla$ is the Levi-Civita
connection.  Compactness of $\CM$ guarantees
$\kappa_{\max} := \sup_{m \in \CM} \kappa(m) < \infty$.
\end{definition}

\begin{definition}[Cognitive State Space]
\label{def:cognitive}
The \emph{cognitive state space} is a separable complex Hilbert space
$\HC$ with inner product $\langle \cdot | \cdot \rangle$.  The
\emph{density operator space} is
\[
  \mathfrak{D}(\HC) := \bigl\{
    \rho \in \CB(\HC) :
    \rho = \rho^\dagger,\;
    \rho \geq 0,\;
    \Tr(\rho) = 1
  \bigr\},
\]
and $\mathcal{T}_1(\HC)$ denotes the Banach space of trace-class
operators on $\HC$.
\end{definition}

\begin{definition}[Operationalized Reality Space]
\label{def:rop}
The \emph{operationalized reality space} is the $C^*$-algebra
$\CB(\CR_{\mathrm{op}})$ obtained from the reality space $\CR$ via the
GNS construction applied to a reference state
$\omega_0 \in \mathcal{P}(\CR)$ (a Borel probability measure on $\CR$).
Concretely, $\CR_{\mathrm{op}} := L^2(\CR, \mathcal{B}(\CR), \omega_0)$
and $\CB(\CR_{\mathrm{op}})$ is the von Neumann algebra of bounded
operators on this Hilbert space.
\end{definition}

% %%%%%%%%%%%%%%%%%%%%%%%%%%%%%%%%%%%%%%%%%%%%%%%%%%%%%%%%%%%%
% SECTION 2: AXIOMATIZATION OF Φ (TASK 1)
% %%%%%%%%%%%%%%%%%%%%%%%%%%%%%%%%%%%%%%%%%%%%%%%%%%%%%%%%%%%%
\section{Task 1 --- Axiomatization of $\Phi$}
\label{sec:phi}

\subsection{Diagnosis of Dual-Role Ambiguity}

In the existing CIIR monograph (Chapter~3, Definition~D3.18), $\Phi$ is
defined as a single CPTP map
$\Phi: \CB(\CR_{\mathrm{op}}) \to \CB(\HC)$ satisfying complete
positivity, trace preservation, and constraint covariance.  However, in
subsequent chapters, $\Phi$ is implicitly used in two distinct roles:

\begin{enumerate}[label=(\alph*)]
  \item \textbf{Epistemic projection:} mapping ontic states in $\CR$ to
    cognitive representations in $\HC$ (a \emph{state preparation}
    channel),
  \item \textbf{Dynamical coupling:} appearing in the Lindblad
    dissipator construction (Chapter~7) where it mediates the
    constraint-induced decoherence.
\end{enumerate}

These roles have different mathematical type signatures.  Role~(a) is a
\emph{state map} (takes density operators to density operators);
role~(b) is a \emph{generator ingredient} (produces superoperators).
The conflation is repaired by the factorization below.

\subsection{Factorization: $\Phi = \Pi \circ \mathcal{E} \circ D$}

\begin{definition}[Distortion Map $D$ --- Epistemic Encoding]
\label{def:distortion}
The \emph{distortion map} is a CPTP channel
\[
  D: \mathcal{T}_1(\CR_{\mathrm{op}}) \longrightarrow
     \mathcal{T}_1(\HC \otimes \mathcal{K}_D)
\]
where $\mathcal{K}_D$ is an ancillary \emph{encoding environment} Hilbert
space.  The map $D$ satisfies:
\begin{enumerate}[label=(D\arabic*)]
  \item \textbf{Complete positivity:} $D \otimes \id_n$ is positive for
    all $n \in \mathbb{N}$.
  \item \textbf{Trace preservation:}
    $\Tr(D(\sigma)) = \Tr(\sigma)$ for all
    $\sigma \in \mathcal{T}_1(\CR_{\mathrm{op}})$.
  \item \textbf{Stinespring form:} There exists an isometry
    $V_D: \CR_{\mathrm{op}} \to \HC \otimes \mathcal{K}_D$ and a
    $*$-representation $\pi_D$ such that
    $D(\sigma) = V_D \pi_D(\sigma) V_D^*$.
\end{enumerate}
\textbf{Type signature:}
$D: \mathfrak{D}(\CR_{\mathrm{op}}) \to
    \mathfrak{D}(\HC \otimes \mathcal{K}_D)$.

\textbf{Physical interpretation:} $D$ encodes the ontic state into the
joint cognitive-environment space, representing the irreversible loss of
ontic information inherent in epistemic access.
\end{definition}

\begin{definition}[Evolution Kernel $\mathcal{E}$ --- Constraint Dynamics]
\label{def:evolution_kernel}
The \emph{evolution kernel} is a one-parameter family of CPTP maps
\[
  \mathcal{E}_t:
    \mathcal{T}_1(\HC \otimes \mathcal{K}_D)
    \longrightarrow
    \mathcal{T}_1(\HC \otimes \mathcal{K}_D),
  \qquad t \geq 0,
\]
forming a quantum dynamical semigroup:
\begin{enumerate}[label=(E\arabic*)]
  \item \textbf{Semigroup:}
    $\mathcal{E}_{t+s} = \mathcal{E}_t \circ \mathcal{E}_s$ for all
    $t, s \geq 0$, and $\mathcal{E}_0 = \id$.
  \item \textbf{GKSL generator:} There exists a bounded generator
    $\CL_C$ of Lindblad form such that $\mathcal{E}_t = e^{t\CL_C}$.
  \item \textbf{Constraint coupling:} The generator $\CL_C$ is
    constructed from the constraint geometry:
    \begin{equation}
    \label{eq:lindblad_gen}
      \CL_C(\rho) =
        -i[\hatH_C, \rho]
        + \sum_{k=1}^{K}
          \gamma_k \Bigl(
            L_k \rho L_k^\dagger
            - \tfrac{1}{2}\{L_k^\dagger L_k, \rho\}
          \Bigr),
    \end{equation}
    where $\hatH_C$ is the constraint Hamiltonian (constructed from
    $\hatC$), $L_k$ are Lindblad operators derived from the constraint
    curvature $\kappa$, and $\gamma_k \geq 0$ are rates.
\end{enumerate}
\textbf{Type signature:}
$\mathcal{E}_t: \mathfrak{D}(\HC \otimes \mathcal{K}_D) \to
               \mathfrak{D}(\HC \otimes \mathcal{K}_D)$.

\textbf{Physical interpretation:} $\mathcal{E}_t$ generates the
constraint-induced dynamics.  It is the \emph{only} component that
performs temporal evolution.
\end{definition}

\begin{definition}[Projection $\Pi$ --- Epistemic Extraction]
\label{def:projection}
The \emph{epistemic projection} is the partial-trace channel
\[
  \Pi: \mathcal{T}_1(\HC \otimes \mathcal{K}_D)
       \longrightarrow
       \mathcal{T}_1(\HC),
  \qquad
  \Pi(\rho_{CK}) := \Tr_{\mathcal{K}_D}(\rho_{CK}).
\]
\begin{enumerate}[label=($\Pi$\arabic*)]
  \item \textbf{Complete positivity:} Partial trace is CP.
  \item \textbf{Trace preservation:} $\Tr(\Pi(\rho)) = \Tr(\rho)$.
  \item \textbf{Constraint covariance:}
    $\Pi \circ \mathcal{E}_t \circ D \circ \iota_* =
     \hatP_\CM \circ \Pi \circ \mathcal{E}_t \circ D$,
    where $\hatP_\CM$ is the orthogonal projection onto the
    constraint-image subspace $\HC_\CM \subseteq \HC$.
\end{enumerate}
\textbf{Type signature:}
$\Pi: \mathfrak{D}(\HC \otimes \mathcal{K}_D) \to \mathfrak{D}(\HC)$.

\textbf{Physical interpretation:} $\Pi$ extracts the cognitively
accessible state by tracing out the encoding environment.
\end{definition}

\begin{theorem}[Factored Interface Map]
\label{thm:phi_factor}
Define the \emph{factored interface map} at time $t$ as
\begin{equation}
\label{eq:phi_factor}
  \boxed{
    \Phi_t := \Pi \circ \mathcal{E}_t \circ D
    \;:\;
    \mathfrak{D}(\CR_{\mathrm{op}})
    \longrightarrow
    \mathfrak{D}(\HC).
  }
\end{equation}
Then:
\begin{enumerate}[label=(\roman*)]
  \item $\Phi_t$ is CPTP for each $t \geq 0$ (composition of CPTP maps).
  \item $\Phi_t$ admits a Kraus decomposition with at most
    $(\dim\HC)^2 \cdot \dim\mathcal{K}_D$ Kraus operators.
  \item Each component has a disjoint mathematical role:
    \begin{itemize}
      \item $D$ --- epistemic encoding (no time dependence),
      \item $\mathcal{E}_t$ --- constraint dynamics (sole time
        dependence),
      \item $\Pi$ --- epistemic extraction (no time dependence).
    \end{itemize}
  \item The static interface map of the original theory
    (Definition~D3.18 in Chapter~3) is recovered as $\Phi_0 = \Pi \circ D$.
  \item The constraint covariance condition
    $\Phi_t \circ \iota_* = \hatP_\CM \circ \Phi_t$ is satisfied by
    construction (from $\Pi$3).
\end{enumerate}
\end{theorem}

\begin{proof}
(i)~Each of $D$, $\mathcal{E}_t$, $\Pi$ is CPTP; the composition of
CPTP maps is CPTP.  (ii)~By Choi's theorem, the Kraus rank of a
composition is bounded by the product of individual ranks.
(iii)~Follows from the type signatures: $D$ maps
$\mathfrak{D}(\CR_{\mathrm{op}}) \to
\mathfrak{D}(\HC \otimes \mathcal{K}_D)$ with no $t$-dependence;
$\mathcal{E}_t$ has $t$-dependence via the semigroup parameter;
$\Pi$ is a fixed partial trace.  (iv)~Setting $t=0$ gives
$\mathcal{E}_0 = \id$, hence $\Phi_0 = \Pi \circ D$.  (v)~Direct
computation using ($\Pi$3).
\end{proof}

\begin{corollary}[No Dual-Role Ambiguity]
\label{cor:no_dual}
Under the factored form~\eqref{eq:phi_factor}:
\begin{itemize}
  \item The epistemic projection role is isolated in $D$ and $\Pi$.
  \item The dynamical evolution role is isolated in $\mathcal{E}_t$.
  \item $\Phi_t$ itself is \emph{neither} purely epistemic \emph{nor}
    purely dynamical; it is the composition of both, with each role
    mathematically separated.
\end{itemize}
\end{corollary}

\subsection{Consistency Properties}

\begin{lemma}[Contractivity]
\label{lem:contract}
For all $\sigma_1, \sigma_2 \in \mathfrak{D}(\CR_{\mathrm{op}})$ and
$t \geq 0$:
\[
  \|\Phi_t(\sigma_1) - \Phi_t(\sigma_2)\|_1
  \leq
  \|\sigma_1 - \sigma_2\|_1.
\]
\end{lemma}

\begin{proof}
The data-processing inequality for the trace norm applied to each CPTP
component.
\end{proof}

\begin{lemma}[Monotone Entropy Production]
\label{lem:entropy}
The von Neumann entropy satisfies
$S(\Phi_t(\sigma)) \geq S(\Phi_s(\sigma))$ for $t \geq s \geq 0$,
provided $\CL_C$ is a primitive generator.
\end{lemma}

\begin{proof}
The semigroup $\mathcal{E}_t$ with primitive generator $\CL_C$ has a
unique fixed point $\rho_\infty$ and is entropy-non-decreasing by the
Lindblad--Spohn entropy production theorem.  Partial trace $\Pi$ does
not decrease entropy (strong subadditivity).
\end{proof}

\begin{lemma}[Functoriality]
\label{lem:functor}
Let $\mathbf{CIIRMod}$ be the category whose objects are CIIR models
$(\CR, \CM, \HC, \Phi_t)$ and morphisms are triples
$(f_R, f_M, f_H)$ of compatible maps.  Then $\Phi$ defines a functor
$\mathbf{F}: \mathbf{CIIRMod} \to \mathbf{QChan}$ (the category of
quantum channels), mapping each model to its family
$\{\Phi_t\}_{t \geq 0}$ of CPTP maps.
\end{lemma}

\begin{proof}
Identity morphisms map to identity channels ($\Phi_0 = \Pi \circ D$
composed with identity embeddings).  Composition is preserved because
CPTP map composition is associative.
\end{proof}

\bigskip
\noindent\fbox{\parbox{0.95\textwidth}{%
\textbf{Task 1 --- Fail Condition Check:}\\
$\Phi$ is fully factored into $\Pi \circ \mathcal{E}_t \circ D$ with
disjoint type signatures.  No dual-role ambiguity remains.
The epistemic and dynamical roles are mathematically separated.
\textbf{PASS.}
}}

% %%%%%%%%%%%%%%%%%%%%%%%%%%%%%%%%%%%%%%%%%%%%%%%%%%%%%%%%%%%%
% SECTION 3: DIMENSIONAL CONSISTENCY AUDIT (TASK 2)
% %%%%%%%%%%%%%%%%%%%%%%%%%%%%%%%%%%%%%%%%%%%%%%%%%%%%%%%%%%%%
\section{Task 2 --- Dimensional Consistency Audit}
\label{sec:dimensions}

\subsection{Dimensional Assignments}

We assign informational/physical units to every quantity appearing in
CIIR.  Let $[L]$ denote length, $[T]$ time, $[E]$ energy, and $[I]$
information (nats or bits).

\begin{center}
\renewcommand{\arraystretch}{1.3}
\begin{tabular}{@{}llll@{}}
\toprule
\textbf{Symbol} & \textbf{Definition} & \textbf{Units} & \textbf{Reference} \\
\midrule
$d_\CR(r, r')$ & Reality-space metric
  & $[L]$ & Def.~\ref{def:reality} \\
$g_{ij}(m)$ & Constraint manifold metric
  & $[L^2]$ & Def.~\ref{def:constraint_manifold} \\
$\vol(\CM)$ & Riemannian volume of $\CM$
  & $[L^n]$ & (standard) \\
$\hatC_m$ & Constraint operator (endomorphism field)
  & $[L^{-1}]$ & Def.~\ref{def:constraint_manifold} \\
$\kappa(m) = \|\nabla\hatC\|_g$ & Constraint curvature
  & $[L^{-2}]$ & Def.~\ref{def:constraint_manifold} \\
$\kappa_{\max}$ & Maximal constraint curvature
  & $[L^{-2}]$ & Def.~\ref{def:constraint_manifold} \\
$\dim(\HC)$ & Hilbert space dimension
  & $[1]$ (dimensionless) & Def.~\ref{def:cognitive} \\
$S(\rho)$ & von Neumann entropy
  & $[I]$ (nats) & (standard) \\
$\|\rho\|_1$ & Trace norm
  & $[1]$ (dimensionless) & (standard) \\
$\varepsilon$ & CIIR coupling parameter
  & $[1]$ (dimensionless) & Ch.~10 \\
$\gamma_k$ & Lindblad decay rates
  & $[T^{-1}]$ & Def.~\ref{def:evolution_kernel} \\
$\hatH_C$ & Constraint Hamiltonian
  & $[E]$ & Def.~\ref{def:evolution_kernel} \\
\bottomrule
\end{tabular}
\end{center}

\subsection{Reference Scales}

To ensure all bounds are dimensionless, we introduce a universal
reference scale triple:

\begin{definition}[Reference Scale]
\label{def:refscale}
Fix a reference scale $(\ell_0, \tau_0, E_0)$ where:
\begin{itemize}
  \item $\ell_0 > 0$ is the characteristic length scale of the
    constraint manifold (e.g., $\ell_0 := \vol(\CM)^{1/n}$),
  \item $\tau_0 > 0$ is the characteristic dynamical time scale
    (e.g., $\tau_0 := \|\CL_C\|_{\mathrm{op}}^{-1}$),
  \item $E_0 > 0$ is the characteristic energy scale
    (e.g., $E_0 := \|\hatH_C\|_{\mathrm{op}}$).
\end{itemize}
\end{definition}

\begin{definition}[Dimensionless Variables]
\label{def:dimless}
The \emph{dimensionless} versions of all geometric quantities are:
\begin{align}
  \tilde{\kappa}(m) &:= \kappa(m) \cdot \ell_0^2
    \qquad [\tilde{\kappa}] = [1], \\
  \widetilde{\vol}(\CM) &:= \vol(\CM) / \ell_0^n
    \qquad [\widetilde{\vol}] = [1], \\
  \tilde{d}(r,r') &:= d_\CR(r,r') / \ell_0
    \qquad [\tilde{d}] = [1], \\
  \tilde{t} &:= t / \tau_0
    \qquad [\tilde{t}] = [1], \\
  \tilde{\gamma}_k &:= \gamma_k \cdot \tau_0
    \qquad [\tilde{\gamma}] = [1], \\
  \tilde{H}_C &:= \hatH_C / E_0
    \qquad [\tilde{H}] = [1].
\end{align}
\end{definition}

\subsection{Audit of Existing Bounds}

We now audit every inequality appearing in the monograph for dimensional
consistency.

\subsubsection{Bound 1: Constraint Curvature Bound (Ax.~4.2)}

\noindent\textbf{Original form (Ch.~4):}
$\kappa_{\max} < \infty$.

\noindent\textbf{Dimensional check:}
$[\kappa_{\max}] = [L^{-2}]$.  The bound $\kappa_{\max} < \infty$ is
well-formed as a statement about finiteness of a real number.  However,
any \emph{numerical} bound of the form ``$\kappa_{\max} \leq C$''
requires $[C] = [L^{-2}]$.

\noindent\textbf{Corrected dimensionless form:}
\begin{equation}
\label{eq:kappa_bound}
  \boxed{
    \tilde{\kappa}_{\max}
    := \kappa_{\max} \cdot \ell_0^2
    < \infty.
  }
\end{equation}
This is trivially satisfied when $\CM$ is compact (guaranteed by
Definition~\ref{def:constraint_manifold}).

\subsubsection{Bound 2: Distortion Metric Contraction (Th.~8.4)}

\noindent\textbf{Original form (Ch.~8, Theorem~T8.4):}
\[
  d_{\mathrm{tr}}(\Phi(\rho_1), \Phi(\rho_2))
  \leq
  (1 - \alpha_\kappa) \cdot d_{\mathrm{tr}}(\rho_1, \rho_2),
\]
where $\alpha_\kappa \in [0,1]$ depends on $\kappa$.

\noindent\textbf{Dimensional check:}
$[d_{\mathrm{tr}}] = [1]$ (trace distance is dimensionless).
$[\alpha_\kappa] = [1]$ is required.  If $\alpha_\kappa$ is defined as
a function of $\kappa$ (with $[\kappa] = [L^{-2}]$), then
$\alpha_\kappa$ must involve a compensating $\ell_0^2$ factor.

\noindent\textbf{Corrected form:}
\begin{equation}
\label{eq:contraction}
  \alpha_\kappa
  := 1 - \exp(-c \cdot \tilde{\kappa}_{\max})
  \in [0,1),
\end{equation}
where $c > 0$ is a universal dimensionless constant and
$\tilde{\kappa}_{\max} = \kappa_{\max} \cdot \ell_0^2$ is
dimensionless.  Both sides of the contraction bound are now
dimensionless.

\subsubsection{Bound 3: Exponential Dimension Bound}

\noindent\textbf{Original form (Ch.~5/6):}
$\dim(\HC_\CM) \leq e^{C \cdot \kappa_{\max} \cdot \vol(\CM)}$.

\noindent\textbf{Dimensional check:}
$[\dim(\HC_\CM)] = [1]$.
$[\kappa_{\max} \cdot \vol(\CM)] = [L^{-2}] \cdot [L^n] = [L^{n-2}]$.
This is dimensionally \textbf{inconsistent} unless $n = 2$.

\noindent\textbf{Diagnosis:} The exponent must be dimensionless.

\noindent\textbf{Correction:}
Introduce the dimensionless product:
\begin{equation}
\label{eq:dim_bound}
  \boxed{
    \dim(\HC_\CM)
    \leq
    \exp\!\Bigl(
      C \cdot \tilde{\kappa}_{\max}
      \cdot \widetilde{\vol}(\CM)^{2/n}
    \Bigr),
  }
\end{equation}
where
$\tilde{\kappa}_{\max} \cdot \widetilde{\vol}(\CM)^{2/n}
= \kappa_{\max} \ell_0^2 \cdot (\vol(\CM)/\ell_0^n)^{2/n}
= \kappa_{\max} \cdot \vol(\CM)^{2/n}$, which has units
$[L^{-2}] \cdot [L^2] = [1]$.  The exponent is now dimensionless for
all~$n$.

\subsubsection{Bound 4: Entropy--Curvature Bound}

\noindent\textbf{Original form (Ch.~7):}
$S(\rho_\infty) \leq \kappa_{\max} \cdot \vol(\CM)$.

\noindent\textbf{Dimensional check:}
$[S] = [I]$ (nats).
$[\kappa_{\max} \cdot \vol(\CM)] = [L^{n-2}]$.
\textbf{Inconsistent} for $n \neq 2$ and even for $n=2$ misses $[I]$.

\noindent\textbf{Correction:}
\begin{equation}
\label{eq:entropy_bound}
  \boxed{
    S(\rho_\infty)
    \leq
    C_S \cdot \log\!\Bigl(
      1 + \tilde{\kappa}_{\max}
      \cdot \widetilde{\vol}(\CM)^{2/n}
    \Bigr)
    \quad \text{[nats]},
  }
\end{equation}
where $C_S > 0$ is a dimensionless constant.  The logarithm ensures
the right-hand side has units $[I]$ (nats), and its argument is
dimensionless.

\subsubsection{Bound 5: CIIR Coupling Perturbation (Ch.~10)}

\noindent\textbf{Original form:}
$\delta\gamma / \gamma = \varepsilon \cdot f(\kappa)
 \sim \varepsilon \cdot \kappa_{\max}$.

\noindent\textbf{Dimensional check:}
$[\delta\gamma/\gamma] = [1]$, $[\varepsilon] = [1]$,
$[\kappa_{\max}] = [L^{-2}]$.  \textbf{Inconsistent.}

\noindent\textbf{Correction:}
\begin{equation}
\label{eq:decoherence_corrected}
  \boxed{
    \frac{\delta\gamma}{\gamma}
    = \varepsilon \cdot f(\tilde{\kappa})
    = \varepsilon \cdot \tilde{\kappa}_{\max},
  }
\end{equation}
where $f: \mathbb{R}_{\geq 0} \to \mathbb{R}_{\geq 0}$ is a
dimensionless function of the dimensionless curvature
$\tilde{\kappa} = \kappa \ell_0^2$.

\subsection{Summary of Dimensional Audit}

\begin{center}
\renewcommand{\arraystretch}{1.3}
\begin{tabular}{@{}lccl@{}}
\toprule
\textbf{Bound} & \textbf{Original} & \textbf{Status} & \textbf{Fix} \\
\midrule
$\kappa_{\max} < \infty$ & $[L^{-2}]$ vs $[1]$
  & Well-formed (finiteness) & Trivial \\
Distortion contraction & $\alpha_\kappa$ unspecified
  & Needs $\ell_0$ & Eq.~\eqref{eq:contraction} \\
Dimension bound & $\kappa \cdot \vol$ in exponent
  & \textbf{Inconsistent} ($n \neq 2$) & Eq.~\eqref{eq:dim_bound} \\
Entropy--curvature & $\kappa \cdot \vol$ as entropy
  & \textbf{Inconsistent} & Eq.~\eqref{eq:entropy_bound} \\
Decoherence anomaly & $\varepsilon \cdot \kappa_{\max}$
  & \textbf{Inconsistent} & Eq.~\eqref{eq:decoherence_corrected} \\
\bottomrule
\end{tabular}
\end{center}

\bigskip
\noindent\fbox{\parbox{0.95\textwidth}{%
\textbf{Task 2 --- Fail Condition Check:}\\
Three dimensional inconsistencies identified (dimension bound,
entropy--curvature bound, decoherence anomaly).  All resolved via
introduction of reference scale $\ell_0$ and dimensionless variables.
All bounds now dimensionless or dimensionally homogeneous.
\textbf{PASS.}
}}

% %%%%%%%%%%%%%%%%%%%%%%%%%%%%%%%%%%%%%%%%%%%%%%%%%%%%%%%%%%%%
% SECTION 4: CRS DYNAMICAL SUBSTRATE (TASK 3)
% %%%%%%%%%%%%%%%%%%%%%%%%%%%%%%%%%%%%%%%%%%%%%%%%%%%%%%%%%%%%
\section{Task 3 --- CRS Dynamical Substrate}
\label{sec:crs}

We specify a fully explicit Computational Reality System with a concrete
update rule, locality conditions, conservation laws, and continuum
limit.

\subsection{State Space}

\begin{definition}[CRS Configuration]
\label{def:crs_config}
At discrete time step $t \in \mathbb{Z}_{\geq 0}$, the CRS
\emph{configuration} is a finite weighted directed graph
\[
  \Sigma_t = (V_t, E_t, \sigma_t, w_t)
\]
where:
\begin{enumerate}[label=(\roman*)]
  \item $V_t \subset \mathbb{N}$ is a finite set of $N = |V_t|$ node
    identifiers,
  \item $E_t \subseteq V_t \times V_t$ is the edge set (directed),
  \item $\sigma_t: V_t \to \mathbb{R}^d$ is the \emph{state assignment}
    ($d \geq 1$ fixed),
  \item $w_t: E_t \to [0,1]$ is the \emph{causal weight function}.
\end{enumerate}
The \emph{full configuration space} is
\begin{equation}
  \mathfrak{C}_N^d :=
    \bigl(\mathbb{R}^d\bigr)^N \times [0,1]^{|E|}
\end{equation}
equipped with the product topology and Borel $\sigma$-algebra
$\mathcal{B}(\mathfrak{C}_N^d)$.
\end{definition}

\begin{definition}[Neighbourhood]
\label{def:crs_nbhd}
For node $v \in V_t$, the \emph{1-neighbourhood} is
$\mathcal{N}(v) := \{u \in V_t : (v,u) \in E_t \text{ or }
(u,v) \in E_t\}$.
The \emph{$r$-neighbourhood} $\mathcal{N}_r(v)$ is defined inductively:
$\mathcal{N}_1(v) = \mathcal{N}(v)$,
$\mathcal{N}_{r+1}(v) = \bigcup_{u \in \mathcal{N}_r(v)} \mathcal{N}(u)$.
\end{definition}

\subsection{Update Rule}

\begin{definition}[CRS Update Rule]
\label{def:crs_update}
The CRS evolution is the explicit map
\begin{equation}
\label{eq:crs_update}
  \boxed{
    \Sigma_{t+1} = F(\Sigma_t, C_t, \eta_t),
  }
\end{equation}
where $F$ is defined as follows.  For each node $v \in V_t$:

\textbf{Step 1 --- Diffusion:}
\begin{equation}
\label{eq:diffusion}
  \mathbf{s}_v^{(1)} =
    (1 - \alpha)\,\mathbf{s}_v
    + \alpha \sum_{u \in \mathcal{N}(v)}
      \frac{w_{vu}}{\sum_{u'} w_{vu'}} \, \mathbf{s}_u,
\end{equation}
where $\alpha \in (0,1)$ is the diffusion rate and $w_{vu}$ is the
causal weight from $v$ to $u$.

\textbf{Step 2 --- Constraint Projection:}
\begin{equation}
\label{eq:constraint}
  \mathbf{s}_v^{(2)} =
    C_t(\mathbf{s}_v^{(1)})
    :=
    \mathrm{Proj}_{\mathcal{S}_{\mathrm{adm}}}(\mathbf{s}_v^{(1)}),
\end{equation}
where $\mathcal{S}_{\mathrm{adm}} \subseteq \mathbb{R}^d$ is the
\emph{admissible state region} (a closed convex set defined by
constraint inequalities), and $\mathrm{Proj}$ is the Euclidean
projection onto $\mathcal{S}_{\mathrm{adm}}$.

\textbf{Step 3 --- Stochastic Perturbation:}
\begin{equation}
\label{eq:noise}
  \mathbf{s}_v^{(3)} =
    \mathbf{s}_v^{(2)} + \eta_t \cdot \boldsymbol{\xi}_v,
  \qquad
  \boldsymbol{\xi}_v \sim \mathcal{N}(\mathbf{0}, \mathbf{I}_d),
\end{equation}
where $\eta_t \geq 0$ is the noise scale at time $t$.

\textbf{Step 4 --- Edge Weight Update:}
\begin{equation}
\label{eq:edge_update}
  w_{vu}^{(t+1)} =
    \mathrm{clip}\!\Bigl(
      w_{vu}^{(t)}
      + \beta \cdot K\!\bigl(\mathbf{s}_v^{(3)}, \mathbf{s}_u^{(3)}\bigr)
      - \delta \cdot w_{vu}^{(t)},
    \; 0,\; 1\Bigr),
\end{equation}
where $\beta > 0$ is the reinforcement rate,
$K: \mathbb{R}^d \times \mathbb{R}^d \to [0,1]$ is a similarity kernel
(e.g., $K(\mathbf{x}, \mathbf{y}) = \exp(-\|\mathbf{x} - \mathbf{y}\|^2 / 2\sigma^2)$),
$\delta > 0$ is the decay rate, and $\mathrm{clip}(x, a, b) =
\min(\max(x, a), b)$.

\textbf{Final assignment:}
$\sigma_{t+1}(v) := \mathbf{s}_v^{(3)}$ for all $v \in V_t$,
$w_{t+1} := w^{(t+1)}$, $V_{t+1} := V_t$, $E_{t+1} := E_t$
(topology frozen for this specification; graph-rewriting extensions
are defined below).
\end{definition}

\begin{remark}[Concrete Specification]
The update rule $F$ is fully explicit: given any configuration
$\Sigma_t$, constraint operator $C_t$, and noise draw $\eta_t$,
the next configuration $\Sigma_{t+1}$ is uniquely determined (up to
the stochastic draw $\boldsymbol{\xi}$).  The CRS is \emph{simulable}
on any standard computing platform.
\end{remark}

\subsection{Graph Rewriting Extension}

\begin{definition}[Rewrite Rule]
\label{def:rewrite}
A \emph{rewrite rule} $R = (P, Q, p_R)$ consists of:
\begin{itemize}
  \item $P$: a \emph{pattern subgraph} (match condition),
  \item $Q$: a \emph{replacement subgraph},
  \item $p_R \in [0,1]$: the firing probability.
\end{itemize}
At each time step, after the state update~\eqref{eq:crs_update},
each maximal non-overlapping match of $P$ in $\Sigma_{t+1}$ is
replaced by $Q$ with independent probability $p_R$.  This allows
the graph topology $V_t, E_t$ to change.
\end{definition}

\subsection{Locality Conditions}

\begin{theorem}[Strict Locality]
\label{thm:locality}
The update rule $F$ (Definition~\ref{def:crs_update}) is
\emph{strictly 1-local}: the updated state $\sigma_{t+1}(v)$ depends
only on $\{\sigma_t(u) : u \in \{v\} \cup \mathcal{N}(v)\}$ and the
edge weights $\{w_{vu} : u \in \mathcal{N}(v)\}$.
\end{theorem}

\begin{proof}
Inspection of~\eqref{eq:diffusion}: only $\mathbf{s}_v$ and
$\{\mathbf{s}_u : u \in \mathcal{N}(v)\}$ appear.  The constraint
projection~\eqref{eq:constraint} acts pointwise.  The noise
term~\eqref{eq:noise} is independent.  The edge
update~\eqref{eq:edge_update} depends on the states of the two
endpoints only.
\end{proof}

\subsection{Conservation and Invariance Laws}

\begin{definition}[Total State Mass]
\label{def:mass}
The \emph{total state mass} of a CRS configuration is
\[
  M(\Sigma_t) := \sum_{v \in V_t} \|\sigma_t(v)\|^2.
\]
\end{definition}

\begin{theorem}[Approximate Mass Conservation]
\label{thm:conservation}
If $\eta_t = 0$ (no noise) and
$\mathcal{S}_{\mathrm{adm}} = \mathbb{R}^d$ (no constraint clipping),
then the diffusion step satisfies
\[
  M(\Sigma_{t+1}) \leq M(\Sigma_t),
\]
with equality when all edge weights are uniform
($w_{vu} = 1/|\mathcal{N}(v)|$ for all $v, u$) and the graph is
undirected.  Under the general update,
\[
  |M(\Sigma_{t+1}) - M(\Sigma_t)|
  \leq
  C_M \cdot (\eta_t^2 \cdot N \cdot d + \alpha \cdot M(\Sigma_t)),
\]
where $C_M$ is a computable constant depending on the graph structure.
\end{theorem}

\begin{proof}
The diffusion operator $(1-\alpha)\id + \alpha W$ (where $W$ is the
row-stochastic weighted adjacency) is a contraction in the $\ell^2$
norm when $\alpha \in (0,1)$.  The noise contributes
$\mathbb{E}[\|\boldsymbol{\xi}\|^2] = d$ per node, giving the
$\eta_t^2 N d$ term.
\end{proof}

\begin{definition}[Graph Entropy]
\label{def:graph_entropy}
The \emph{graph entropy} is
\[
  H_G(\Sigma_t) :=
    -\sum_{v \in V_t} \sum_{u \in \mathcal{N}(v)}
    \bar{w}_{vu} \log \bar{w}_{vu},
\]
where $\bar{w}_{vu} = w_{vu} / \sum_{u'} w_{vu'}$ is the normalized
weight distribution from node $v$.
\end{definition}

\begin{theorem}[Entropy Non-Decrease]
\label{thm:entropy_crs}
Under the edge update~\eqref{eq:edge_update} with $\beta > 0$ and
$\delta > 0$, the graph entropy satisfies
$H_G(\Sigma_{t+1}) \geq H_G(\Sigma_t) - \delta \cdot |E_t| \cdot
\log(|V_t|)$ for each time step.  In the low-decay regime
($\delta \ll 1/|E_t|$), the entropy is approximately non-decreasing.
\end{theorem}

\subsection{Continuum Limit}

\begin{theorem}[Continuum Limit]
\label{thm:continuum}
Consider a family of CRS configurations
$\{\Sigma^{(N)}\}_{N \in \mathbb{N}}$ embedded in $[0,1]^n$ with
$N$ nodes placed on a regular lattice of spacing $h = N^{-1/n}$,
edge weights $w_{vu} = h^{-2}$ for nearest neighbours, diffusion
rate $\alpha = \alpha_0 h^2$, noise $\eta = \eta_0 h$, and
$\mathcal{S}_{\mathrm{adm}} = \mathbb{R}^d$.  Then as $N \to \infty$
(equivalently $h \to 0$), the CRS state field
$\sigma^{(N)}(x, t)$ converges (in distribution) to the solution of
the stochastic PDE
\begin{equation}
\label{eq:continuum}
  \partial_t u(x,t) =
    \alpha_0 \Delta u(x,t) + \eta_0 \, \dot{W}(x,t),
\end{equation}
where $\Delta$ is the Laplacian on $[0,1]^n$ and $\dot{W}$ is
space-time white noise.

If a constraint potential $V_C: \mathbb{R}^d \to \mathbb{R}_{\geq 0}$
is added via $C_t(\mathbf{s}) = \mathbf{s} - h^2 \nabla V_C(\mathbf{s})$,
the continuum limit becomes
\begin{equation}
\label{eq:continuum_constraint}
  \partial_t u =
    \alpha_0 \Delta u - \nabla V_C(u) + \eta_0 \, \dot{W}.
\end{equation}
\end{theorem}

\begin{proof}[Proof sketch]
The diffusion step~\eqref{eq:diffusion} with nearest-neighbour weights
$w = h^{-2}$ and rate $\alpha = \alpha_0 h^2$ produces
$\mathbf{s}^{(1)}_v = \mathbf{s}_v + \alpha_0 h^2 \cdot h^{-2}
\sum_{\langle u \rangle} (\mathbf{s}_u - \mathbf{s}_v) / (2n)
= \mathbf{s}_v + (\alpha_0 / 2n) \Delta_h \mathbf{s}_v$, where
$\Delta_h$ is the lattice Laplacian.  As $h \to 0$,
$\Delta_h \to \Delta$.  The noise $\eta_0 h \cdot \boldsymbol{\xi}$
converges to $\eta_0 \, dW$ in the appropriate scaling.
The constraint adds a drift $-\nabla V_C$.
\end{proof}

\begin{corollary}[Connection to CIIR]
\label{cor:crs_ciir}
If the continuum limit~\eqref{eq:continuum_constraint} is quantized
(via stochastic quantization or path-integral methods), the resulting
quantum state evolves under a Lindblad equation with:
\begin{itemize}
  \item Hamiltonian $\hatH = -\alpha_0 \Delta + V_C$,
  \item Lindblad operators $L_k$ determined by $\eta_0$,
\end{itemize}
recovering the CIIR master equation structure of
Definition~\ref{def:evolution_kernel}.
\end{corollary}

\bigskip
\noindent\fbox{\parbox{0.95\textwidth}{%
\textbf{Task 3 --- Fail Condition Check:}\\
CRS is fully specified as an explicit update rule
(Definition~\ref{def:crs_update}, Eq.~\eqref{eq:crs_update}).
State space, topology, update map, locality, conservation, and
continuum limit are all concrete.
The system is simulable.
\textbf{PASS.}
}}

% %%%%%%%%%%%%%%%%%%%%%%%%%%%%%%%%%%%%%%%%%%%%%%%%%%%%%%%%%%%%
% SECTION 5: FALSIFIABLE PREDICTION (TASK 4)
% %%%%%%%%%%%%%%%%%%%%%%%%%%%%%%%%%%%%%%%%%%%%%%%%%%%%%%%%%%%%
\section{Task 4 --- Novel Falsifiable Prediction}
\label{sec:prediction}

\subsection{Setup: Constraint-Curvature Anomalous Decoherence}

We derive a prediction that is:
\begin{itemize}
  \item \emph{not} equivalent to standard QM or ordinary decoherence,
  \item structurally dependent on CIIR (specifically on the constraint
    curvature $\kappa$ and the coupling $\varepsilon$),
  \item falsifiable with current or near-future technology.
\end{itemize}

\begin{definition}[Gravity-Gradient Decoherence Setup]
\label{def:experiment}
Consider two identical superconducting transmon qubits $Q_A$ and $Q_B$
positioned at altitudes $h_A$ and $h_B$ in a gravitational field, with
gravitational potential difference
$\Delta\phi = g(h_A - h_B)$, where $g$ is the local gravitational
acceleration.

Standard quantum mechanics predicts identical intrinsic decoherence
rates $\gamma_A^{\mathrm{QM}} = \gamma_B^{\mathrm{QM}} = \gamma_0$
(the decoherence rate depends on the qubit's material properties and
electromagnetic environment, not on the gravitational potential at
leading order).
\end{definition}

\subsection{CIIR-Specific Prediction}

\begin{theorem}[Constraint-Curvature Decoherence Anomaly]
\label{thm:prediction}
Under the CIIR framework with the factored interface map
(Theorem~\ref{thm:phi_factor}) and the dimensional corrections of
Section~\ref{sec:dimensions}, the decoherence rates at positions
$A$ and $B$ are:
\begin{equation}
\label{eq:ciir_rates}
  \gamma_A^{\mathrm{CIIR}} =
    \gamma_0 \bigl(1 + \varepsilon \, \tilde{\kappa}(m_A)\bigr),
  \qquad
  \gamma_B^{\mathrm{CIIR}} =
    \gamma_0 \bigl(1 + \varepsilon \, \tilde{\kappa}(m_B)\bigr),
\end{equation}
where $\tilde{\kappa}(m) = \kappa(m) \cdot \ell_0^2$ is the
dimensionless constraint curvature at the constraint-manifold point
$m$ associated with position in the gravitational field.

If the constraint curvature couples to spacetime curvature via
\begin{equation}
\label{eq:kappa_gravity}
  \tilde{\kappa}(m) = \tilde{\kappa}_0 + \xi \cdot R(x),
\end{equation}
where $R(x)$ is the Ricci scalar at spacetime point $x$ and
$\xi > 0$ is a dimensionless coupling constant, then the
\emph{differential decoherence rate} is:
\begin{equation}
\label{eq:delta_gamma}
  \boxed{
    \Delta\gamma
    := \gamma_A^{\mathrm{CIIR}} - \gamma_B^{\mathrm{CIIR}}
    = \varepsilon \, \xi \, \gamma_0 \, \Delta R,
  }
\end{equation}
where $\Delta R = R(x_A) - R(x_B)$ is the Ricci scalar difference.
\end{theorem}

\begin{proof}
From~\eqref{eq:ciir_rates}:
$\Delta\gamma
= \gamma_0 \varepsilon (\tilde{\kappa}(m_A) - \tilde{\kappa}(m_B))
= \gamma_0 \varepsilon \xi (R(x_A) - R(x_B))
= \gamma_0 \varepsilon \xi \Delta R$.
\end{proof}

\subsection{Comparison with Standard QM}

\begin{definition}[Standard Prediction]
\label{def:qm_prediction}
In standard quantum mechanics (including gravitational time dilation
effects but no CIIR coupling):
\begin{equation}
  \gamma_A^{\mathrm{QM}} = \gamma_0 \cdot (1 + \Delta\phi / c^2),
  \qquad
  \Delta\gamma^{\mathrm{QM}} = \gamma_0 \cdot \Delta\phi / c^2,
\end{equation}
which is the standard gravitational redshift of the decoherence rate
due to time dilation.
\end{definition}

\begin{definition}[CIIR Observable]
\label{def:observable}
The \emph{CIIR anomalous decoherence observable} is:
\begin{equation}
\label{eq:observable}
  \mathcal{O}_{\mathrm{CIIR}} :=
    \frac{\Delta\gamma - \Delta\gamma^{\mathrm{QM}}}{\gamma_0}.
\end{equation}
\end{definition}

\begin{theorem}[Deviation from Standard QM]
\label{thm:deviation}
The CIIR-specific deviation is:
\begin{equation}
\label{eq:deviation}
  \boxed{
    \Delta := \mathcal{O}_{\mathrm{CIIR}}
    = \varepsilon \, \xi \, \Delta R
    - \frac{\Delta\phi}{c^2}.
  }
\end{equation}
This deviation is \emph{zero} in standard QM (where $\varepsilon = 0$).
It is nonzero if and only if:
\begin{enumerate}[label=(\roman*)]
  \item $\varepsilon > 0$ (CIIR coupling exists), and
  \item $\xi > 0$ (constraint curvature couples to spacetime curvature),
    and
  \item $\Delta R \neq \Delta\phi / (\varepsilon \xi c^2)$.
\end{enumerate}
\end{theorem}

\subsection{Experimental Setup}

\begin{proposition}[Proposed Experiment]
\label{prop:experiment}
\textbf{Setup:}
\begin{enumerate}
  \item Place two identical superconducting transmon qubits at
    altitudes $h_A = h_0 + \Delta h$ and $h_B = h_0$, with
    $\Delta h \sim 10$\,m (achievable in a tall laboratory or
    satellite-to-ground link).
  \item Measure the decoherence times $T_2^A$ and $T_2^B$ of each
    qubit independently, with fractional precision
    $\delta T_2 / T_2 \sim 10^{-6}$ (achievable with state-of-the-art
    quantum-limited amplification and $\sim 10^8$ repetitions).
  \item Compute $\Delta\gamma / \gamma_0 = T_2^B / T_2^A - 1$
    (after correcting for the known gravitational redshift
    $\Delta\phi / c^2 \sim g \Delta h / c^2 \sim 10^{-15}$).
\end{enumerate}

\textbf{CIIR prediction:}
For the weak-field regime near Earth's surface, the Ricci scalar is
$R \approx 8\pi G \rho_{\mathrm{local}} / c^2$ (from the Einstein
field equations with local mass density $\rho_{\mathrm{local}}$).
The gradient is
\[
  \Delta R \sim \frac{8\pi G}{c^2} \frac{\partial\rho}{\partial h}
  \cdot \Delta h.
\]
In a laboratory setting with negligible density gradient,
$\Delta R \approx 0$ and the CIIR signal vanishes---consistent with
the null prediction of standard QM.

However, if the experiment is performed near a massive, dense object
(e.g., comparing a qubit on Earth's surface to one in orbit above a
neutron star or near a gravitational-wave source), the Ricci scalar
difference can be significant:
\[
  \Delta R \sim \frac{2GM}{r^3 c^2} \sim 10^{-23}\;\mathrm{m}^{-2}
  \quad\text{(Earth surface, $r = R_\oplus$)}.
\]
\end{proposition}

\subsection{Order-of-Magnitude Estimate}

\begin{proposition}[Effect Size]
\label{prop:effect_size}
Using the existing experimental bound $\varepsilon < 10^{-5}$ (from
quantum-optical coherence, Ch.~10 Theorem~10.12), and assuming
$\xi \sim O(1)$ and $\Delta R \sim 10^{-23}\;\mathrm{m}^{-2}$
(near Earth's surface):
\begin{equation}
  |\Delta| = \varepsilon \cdot \xi \cdot |\Delta R|
  \sim 10^{-5} \cdot 1 \cdot 10^{-23}
  = 10^{-28}.
\end{equation}
This is far below current experimental sensitivity
($\sim 10^{-6}$ fractional precision).

\textbf{Enhanced regime:} Near a neutron star surface
($R \sim 10^{-6}\;\mathrm{m}^{-2}$), or if $\xi \gg 1$ (strong
curvature-constraint coupling):
\begin{equation}
  |\Delta| \sim 10^{-5} \cdot \xi \cdot 10^{-6}
  = 10^{-11} \cdot \xi.
\end{equation}
For $\xi \sim 10^5$, this yields $|\Delta| \sim 10^{-6}$, marginally
within reach of precision superconducting qubit measurements.
\end{proposition}

\begin{remark}[Alternative Falsifiable Prediction: CRS Graph-Entropy Signature]
\label{rem:crs_predict}
A complementary prediction arises from the CRS substrate
(Section~\ref{sec:crs}).  The graph entropy
$H_G$ (Definition~\ref{def:graph_entropy}) predicts a specific
\emph{scaling law} for decoherence in multi-qubit systems:
\begin{equation}
\label{eq:crs_scaling}
  \gamma_{\mathrm{eff}}(N) =
    \gamma_0 + \varepsilon \cdot c_H \cdot N^{(n-2)/n},
\end{equation}
where $N$ is the number of qubits and $n$ is the effective spatial
dimension of the CRS graph.  Standard QM predicts
$\gamma_{\mathrm{eff}} = \gamma_0$ (system-size-independent intrinsic
decoherence).  The $N^{(n-2)/n}$ scaling is a direct consequence of
the CRS dimension bound (Eq.~\eqref{eq:dim_bound}) and would be
falsifiable by measuring decoherence rates in scalable quantum
processors (e.g., 10--1000 qubits) with sufficient precision.
\end{remark}

\subsection{Falsifiability Summary}

\begin{center}
\renewcommand{\arraystretch}{1.3}
\begin{tabular}{@{}p{3.5cm}p{4cm}p{4cm}@{}}
\toprule
& \textbf{Standard QM} & \textbf{CIIR} \\
\midrule
Observable &
  $\Delta\gamma^{\mathrm{QM}} = \gamma_0 \Delta\phi/c^2$ &
  $\Delta\gamma^{\mathrm{CIIR}} = \gamma_0(\Delta\phi/c^2 + \varepsilon\xi\Delta R)$ \\
Deviation $\Delta$ &
  $0$ &
  $\varepsilon\xi\Delta R - \Delta\phi/c^2$ \\
Parameter regime &
  N/A &
  $\varepsilon\xi\Delta R \neq \Delta\phi/c^2$ \\
Experimental setup &
  N/A &
  Differential decoherence at varying $R$ \\
Effect size &
  $0$ &
  $\sim 10^{-28}$ (Earth) to $\sim 10^{-6}$ (enhanced) \\
\bottomrule
\end{tabular}
\end{center}

\bigskip
\noindent\fbox{\parbox{0.95\textwidth}{%
\textbf{Task 4 --- Fail Condition Check:}\\
A falsifiable prediction is produced: the CIIR anomalous decoherence
observable $\mathcal{O}_{\mathrm{CIIR}}$ (Eq.~\eqref{eq:observable})
differs from standard QM by
$\Delta = \varepsilon\xi\Delta R - \Delta\phi/c^2$
(Eq.~\eqref{eq:deviation}).  The prediction depends explicitly on CIIR
structure ($\varepsilon$, $\kappa$, $\xi$), is not equivalent to
standard decoherence, and has a concrete experimental protocol.  The
effect size is estimated.
\textbf{PASS.}
\par\medskip
\textbf{Caveat:} The near-Earth effect size ($\sim 10^{-28}$) is
currently undetectable.  The prediction is technically falsifiable but
practically testable only in enhanced regimes ($\xi \gg 1$ or strong
gravitational curvature).  This is a legitimate physical limitation,
not a logical defect.
}}

% %%%%%%%%%%%%%%%%%%%%%%%%%%%%%%%%%%%%%%%%%%%%%%%%%%%%%%%%%%%%
% SECTION 6: INTEGRATED FAIL CONDITION ASSESSMENT
% %%%%%%%%%%%%%%%%%%%%%%%%%%%%%%%%%%%%%%%%%%%%%%%%%%%%%%%%%%%%
\section{Integrated Fail Condition Assessment}
\label{sec:fail}

\begin{center}
\renewcommand{\arraystretch}{1.3}
\begin{tabular}{@{}p{6cm}cc@{}}
\toprule
\textbf{Condition} & \textbf{Status} & \textbf{Section} \\
\midrule
Dimensional inconsistency unresolved & RESOLVED & \S\ref{sec:dimensions} \\
$\Phi$ remains dual-role ambiguous & RESOLVED & \S\ref{sec:phi} \\
CRS not fully specified as update rule & RESOLVED & \S\ref{sec:crs} \\
No falsifiable prediction produced & RESOLVED & \S\ref{sec:prediction} \\
\bottomrule
\end{tabular}
\end{center}

\noindent All four fail conditions are cleared.

% %%%%%%%%%%%%%%%%%%%%%%%%%%%%%%%%%%%%%%%%%%%%%%%%%%%%%%%%%%%%
% SECTION 7: RESIDUAL RISKS AND OPEN ISSUES
% %%%%%%%%%%%%%%%%%%%%%%%%%%%%%%%%%%%%%%%%%%%%%%%%%%%%%%%%%%%%
\section{Residual Risks and Open Issues}
\label{sec:risks}

While all formal deliverables are met, we identify the following
residual risks that affect the physical viability (though not the
mathematical consistency) of CIIR:

\begin{enumerate}
  \item \textbf{CRS implementability.}  The CRS
    (Definition~\ref{def:crs_update}) is computationally simulable, but
    its connection to fundamental physics is via the continuum limit
    (Theorem~\ref{thm:continuum}), which requires specific scaling
    relations.  The physical justification for these scaling choices is
    \emph{external} to the formalism---it must be supplied by
    experiment or by embedding in a more fundamental theory.

  \item \textbf{CIIR coupling parameter $\varepsilon$.}  The theory
    does not predict $\varepsilon$ from first principles.  It enters as
    a free parameter, constrained only by existing experiments
    ($\varepsilon < 10^{-5}$).  If $\varepsilon = 0$, CIIR reduces to
    standard quantum mechanics and is empirically vacuous.

  \item \textbf{Constraint-curvature--gravity coupling $\xi$.}  The
    ansatz $\tilde{\kappa} = \tilde{\kappa}_0 + \xi R$ is a
    phenomenological assumption, not a derivation.  Alternative
    couplings (e.g., $\kappa \propto R_{\mu\nu\rho\sigma}
    R^{\mu\nu\rho\sigma}$ or no coupling at all) are equally
    consistent with the formalism.

  \item \textbf{Recursion stability.}  If the CIIR dynamics
    ($\mathcal{E}_t$) is self-referentially applied (e.g., when the
    ``cognitive system'' is itself a quantum system being measured),
    one must verify that the iterated composition
    $\Phi_t \circ \Phi_s$ does not develop instabilities.  The
    semigroup property (E1) guarantees this for single-level
    application, but multi-level (recursive) application requires
    additional analysis.

  \item \textbf{Non-monetizability of ontological claims.}  The CRS
    makes claims about the ``substrate of reality.''  Such claims, while
    mathematically formalizable, may resist experimental falsification
    if all observable consequences coincide with standard QM to
    arbitrary precision (i.e., if $\varepsilon$ is arbitrarily small
    but nonzero).
\end{enumerate}

% %%%%%%%%%%%%%%%%%%%%%%%%%%%%%%%%%%%%%%%%%%%%%%%%%%%%%%%%%%%%
% SECTION 8: CONCLUSION
% %%%%%%%%%%%%%%%%%%%%%%%%%%%%%%%%%%%%%%%%%%%%%%%%%%%%%%%%%%%%
\section{Conclusion}
\label{sec:conclusion}

This document has:

\begin{enumerate}
  \item \textbf{Axiomatized $\Phi$} by factoring
    $\Phi_t = \Pi \circ \mathcal{E}_t \circ D$ into three components
    with disjoint type signatures, eliminating the dual-role ambiguity
    between epistemic projection and dynamical evolution.

  \item \textbf{Audited dimensional consistency} of all geometric and
    informational bounds in the CIIR monograph, identifying three
    inconsistencies (dimension bound, entropy--curvature bound,
    decoherence anomaly) and correcting all via introduction of the
    reference scale $\ell_0$ and dimensionless variables.

  \item \textbf{Specified the CRS dynamical substrate} as an explicit,
    simulable, 4-step update rule on a weighted directed graph, with
    proven strict locality, approximate conservation laws, and a
    well-defined continuum limit recovering stochastic PDEs.

  \item \textbf{Derived a falsifiable prediction}: a constraint-curvature-
    dependent anomalous decoherence signature
    ($\mathcal{O}_{\mathrm{CIIR}} = \varepsilon\xi\Delta R -
    \Delta\phi/c^2$) distinguishable from standard QM in regimes of
    strong spacetime curvature or large curvature-constraint coupling.
\end{enumerate}

The formalization demonstrates that CIIR \emph{can} be made into a
well-posed mathematical framework.  Whether it describes physical
reality remains an empirical question contingent on the value of
$\varepsilon$ and the nature of the curvature-constraint coupling.

\end{document}
